\subsection{$\SL(n,k)$-invariance}

Let $A=(\a{i, j})$ be the $n\times (n+1)$ matrix of variables. For a matrix $B\in \SL(n,k)$, consider the linear change of variables from $A$ to $BA$. This defines an action of $\SL(n,k)$ on the polynomial ring $k[\a{i, j}]$. The first fundamental theorem of invariant theory for $\SL(n,k)$ states that if $k$ is an infinite field of any characteristic,  then the $k$-algebra of invariants under this action is generated by $X_0, X_1, \dots, X_n$. When  the field $k$ is finite, the same result holds if we consider absolute invariants. These are polynomials in $k[\a{i,j}]$ that are invariant under the action of $\SL(n,\bar{k})$, where $\bar{k}$ is the algebraic closure of $k$. The first fundamental theorem states that absolute invariants are again polynomials in $X_0, X_1, \dots, X_n$ with coefficients in $k$. (See \cite{Procesi2007}, Theorem~13.5.5 and the discussion of absolute invariants in  Section 13.6.1.)

\begin{lemma} \label{lem-invariance}
Let $I$ and $J$ be as in Theorem~\ref{thm-simplest}. Then the polynomial $\partial_I X_J \in k[\a{i, j}]$ is $\SL(n,k)$-invariant for any field $k$ of characteristic $2$. In particular, $\partial_I X_J \in \FF_2[\a{i, j}]$ is an absolute $\SL(n, \FF_2)$-invariant.
\end{lemma}

\begin{proof}
The group $\SL(n,k)$ is generated by elementary matrices. An elementary matrix acts on the matrix of variables by adding a constant $c$ times row $r$ to row $s$. It suffices to prove invariance under this change of variables. 

We may assume without loss of generality that $r=2$ and $s=1$. Consider the new variables
\[ a'_{i,j} = \begin{cases} 
\a{i, j} + c \a{2, j} & \text{if $i=1$,} \\ 
\a{i, j} & \text{otherwise.}
\end{cases}
\]
For a polynomial $f(a) = f(\a{i, j})$,  let us denote by $f(a')$ the result of substituting $\ap{i,j}$ in $\a{i, j}$. Similarly, let us write $\partial_I(a')$ for the partial derivative where we replace $\partial_{\a{i, j}}$ with $\partial_{a'_{i,j}}$. We need to prove that 
\[ (\partial_I X_J)(a') = (\partial_I X_J) (a). \]
We claim that if $\partial_I = \partial_{\a{1, i_1}}  \partial_{\a{2, i_2}} \partial_{\a{3, i_3}}\cdots \partial_{\a{n, i_n}}$, then 
\begin{equation} (\partial_I X_J)(a') = (\partial_I X_J)(a) - c (\partial_{\a{1, i_1}}  \partial_{\a{1, i_2}} \partial_{\a{3, i_3}} \cdots \partial_{\a{n, i_n}} X_J)(a). \label{eq-chain} \end{equation}
The next lemma shows that $\partial_{\a{1, i_1}}  \partial_{\a{1, i_2}} X_J = 0$, hence the second summand vanishes.

For any polynomial $f(\a{i, j})$, by replacing all symbols $a$ with $a'$, we have  
\[ (\partial_I f)(a')= \partial_I(a') f(a').\]
If $f$ is $\SL(n,k)$-invariant, then 
\[ f(a') = f(a) = f(a'_{1,j} - c a'_{2,j}, a'_{2,j},\ldots, a'_{n,j}).\] 
Let us now prove Equation~\eqref{eq-chain}. Since $X_J$ is $\SL(n,k)$-invariant, 
\[  (\partial_I X_J)(a') = \partial_I(a') X_J(a') = \partial_{\ap{1, i_1}}  \partial_{\ap{2, i_2}} \cdots \partial_{\ap{n, i_n}} X_J (a'_{1,j} - c a'_{2,j}, a'_{2,j},\ldots, a'_{n,j}).\]
Using the chain rule, this derivative is
\[ \big( \partial_I X_J - c \partial_{\a{1, i_1}}  \partial_{\a{1, i_2}}\partial_{\a{3, i_3}} \cdots \partial_{\a{n, i_n}} X_J \big) (a'_{1,j} - c a'_{2,j}, a'_{2,j},\ldots, a'_{n,j}).\]
Changing back to the variables $\a{i,j}$ gives the right hand side of  \eqref{eq-chain}.
\end{proof}

\begin{lemma} \label{lem-2deriv}
If $|J|$ is odd then $\partial_{\a{r, i_1}}  \partial_{\a{r, i_2}} X_J = 0$  for any $r, i_1, i_2$.
\end{lemma}

\begin{proof}
  It is enough to consider the case where $J$ contains no repeating indices, since any square factors of $X_J$ can be factored out of the partial derivatives. Under a suitable relabelling of rows and columns of $A$, we can assume that $r=1$, $i_1 = 1$ and $i_2 = 2$, so that the derivative under consideration is $\partial_{\a{1, 1}} \partial_{\a{1, 2}}  X_J$.
  
 Let us denote by $Y_{i,j} = Y_{j,i}$ the determinant of the matrix $A$ with its first row and columns $i,j$ removed. Then
 \[ \partial_{\a{1, i}} X_j = Y_{i,j}.\]
 The polynomials $X_i$ and $Y_{i,j}$ satisfy the following relations. For any distinct indices $i,j,p,q$ in increasing order
 \begin{equation} \label{eq-Plucker}
 Y_{i,j} Y_{p,q} - Y_{i,p} Y_{j,q} + Y_{i,q} Y_{j,p} = 0,
 \end{equation}
 and for any distinct indices $i,j,p$ in increasing order
 \begin{equation} \label{eq-flag}
 Y_{i,j} X_p - Y_{i,p} X_{j} + Y_{j,p} X_{i} = 0.
 \end{equation}
 These equations hold in any characteristic. In characteristic $2$ the signs in the equations are not important.
 The equations come from the Pl\"ucker embedding of the Grassmannian. Rows $2,3,\ldots,n$ of the  matrix $A$ span an $(n-1)$-plane in the $(n+1)$-space and hence define a point in the Grassmannian $\Gr(n-1, n+1)$. The polynomials $Y_{i,j}$ are the Pl\"ucker coordinates on this Grassmannian. These coordinates satisfy the Pl\"ucker relations in Equation~\eqref{eq-Plucker}. Similarly, the $n$ rows of the matrix $A$ define a point in $\Gr(n,n+1)$ with coordinates $X_i$. The relations in Equation~\eqref{eq-flag}  state that the $(n-1)$-plane with coordinates $Y_{i,j}$ lies in the $n$-plane with coordinates $X_i$.
 
Using the product rule we have
\[ \partial_{\a{1, 1}}  \partial_{\a{1, 2}} X_J = \sum \frac{X_J}{X_{i}X_{j}} Y_{1,i} Y_{2,j}.\]
Here the sum runs over all vectors $(i,j)$ where $i,j$ are distinct entries of $J$ such that $i\neq 1$ and $j\neq 2$. We claim that this sum is equal to 
\[ \sum \frac{X_J}{X_{i}X_{j}} Y_{1,2} Y_{i,j},\]
  where the sum now runs over all two element subsets $\{i,j\}$ of entries in $J$. To see this, first consider the case where $i$ and $j$ are both distinct from $1$ and $2$. In this case we apply the Pl\"ucker relation to get 
 \[  Y_{1,i} Y_{2,j} + Y_{2,i} Y_{1,j}= Y_{1,2} Y_{i,j}.\]
 The cases where $i=2$ or $j=1$ are simpler and do not require any relation.
 
 We are now reduced to proving that 
 \[  \sum_{\{i,j\}} \frac{X_J}{X_{i}X_{j}} Y_{1,2} Y_{i,j} = X_J Y_{1,2} \sum_{\{i,j\}} \frac{Y_{i,j}}{X_{i}X_{j}}  = 0.\]
Using Equation~\eqref{eq-flag} we have for any distinct $i,j, p$
\[ \frac{Y_{i,j}}{X_{i}X_{j}} + \frac{Y_{i,p}}{X_{i}X_{p}} +\frac{Y_{j,p}}{X_{j}X_{p}} = 0.\]
Now consider all three element subsets $\{i,j,p\}$ of entries in $J$. Then 
\[ \sum_{\{i,j,p\}} \left( \frac{Y_{i,j}}{X_{i}X_{j}} + \frac{Y_{i,p}}{X_{i}X_{p}} +\frac{Y_{j,p}}{X_{j}X_{p}} \right) = 0.\]
Since every pair $\{i,j\}$ occurs in an odd number of triples $\{i,j,p\}$, this sum is equal to 
\[ \sum_{\{i,j\}} \frac{Y_{i,j}}{X_{i}X_{j}}. \qedhere\]
\end{proof}
